\documentclass[10pt,a4paper]{amsart}
\usepackage[margin=19mm]{geometry}
\usepackage{amsmath,amssymb,mathtools}
\usepackage[scr=rsfs,cal=euler]{mathalfa}
\usepackage{xfrac}
\usepackage[unicode,hidelinks,bookmarks=false]{hyperref}
\newcommand{\ie}{i.e.\ }
\newcommand{\cf}{cf.\ }
\newcommand{\resp}{respectively\ }
\newcommand{\Subsection}{\subsection}
\providecommand{\nobreakdash}{\nobreak\hbox{-}}
\setlength{\emergencystretch}{2em}
\multlinegap0pt
\usepackage{amssymb}
\usepackage{mathtools}
\newtheorem{lemma}{Lemma}
\newtheorem{assumption}{Assumption}
\newcommand{\K}{\mathcal{K}}
\newcommand{\bfI}{\boldsymbol{\mathrm{I}}}
\newcommand{\bff}{\boldsymbol{f}}
\newcommand{\bftau}{\boldsymbol{\tau}}
\newcommand{\bfu}{\boldsymbol{u}}
\newcommand{\bfv}{\boldsymbol{v}}
\newcommand{\bfw}{\boldsymbol{w}}
\newcommand{\bfxi}{\boldsymbol{\xi}}
\newcommand{\dd}{\:\mathrm{d}}
\DeclareMathOperator{\dev}{dev}
\DeclareMathOperator{\dist}{dist}
\renewcommand\div{\operatorname{\mathrm{div}}}
\DeclareMathOperator{\supp}{supp}
\title{Weak Solutions for the Unregularised Hibler Momentum Equation with Degenerate Coefficients}
\author{Henrik Schneider}
\date{Source: arXiv:2609.12925v1 (CC BY 4.0)}
\begin{document}
\maketitle

\noindent\emph{Source and licence.} Weak Solutions for the Unregularised Hibler Momentum Equation with Degenerate Coefficients. Version arXiv:2609.12925v1 (CC BY 4.0), distributed by arXiv under the Creative Commons Attribution 4.0 International licence, http://creativecommons.org/licenses/by/4.0/, as recorded in the arXiv licence metadata for this exact version and shown on the abstract page https://arxiv.org/abs/2609.12925v1. Full licence notice: \texttt{licenses/arxiv-2609.12925-CC-BY-4.0.txt}.\par\smallskip

\noindent\textit{Editorial scope.} Counted unit: the article's lemma lem:a3 (source0.tex lines 1830--1836). The article's complete original proof of it is delivered with it (source0.tex lines 1837--1953, one \texttt{proof} environment of 117 lines). No mathematics is written, repaired or re-derived: the statement and the proof are the article's own lines, in the article's own notation, and no retained line is substituted or re-flowed.  Retained uncounted as the source-local context the proof needs: lines 205--212; lines 230--238; lines 261--280; lines 1807--1809; lines 1815--1818.  Nothing was omitted from any retained span, so the package carries no omission record.  No retained line contains a citation, so the wrapper carries no bibliography and no bibliography entry.  No retained line contains a figure environment, an image-inclusion command or drawing code, so no image is delivered and the package declares no asset dependency.  Editorial formatting only, in addition: an AMS \texttt{amsart} preamble that restates the article's own declarations and macros used by the retained lines, namely the environments \texttt{lemma}, \texttt{assumption} on separate counters, together with the standard packages those lines need.  The retained environments print in this document as Lemma 1, Assumption 1 in order of appearance, whereas the article prints the same items with the numbering of its own full document.  The complete native source of the article as distributed in the arXiv e-print for this exact version is bundled unchanged as \texttt{sources/210163/source0.tex}.
\begin{lemma} \label{lema:prop_K}
    For $P\geq 0$ the following statements hold:
    \begin{enumerate}
        \item $\K(P) = P \K(1)$, in particular $P\mapsto \K(P)$ is monotone, i.e.\ from $0\leq P' \leq P$ follows $\K(P') \subseteq \K(P)$, and $0\in \K(P)$ for all $P\geq 0$.
        \item Boundedness: $|\bftau |_F \leq c_\K P$ for all $\bftau \in \K(P)$ with $c_\K = \sqrt{2+1/(2e^2)}$.
        \item For $P>0$, $\K(P)$, understood in the three-dimensional space $(\tau_I, \dev\bftau) \in \mathbb R \times \mathbb R^2$, is an ellipsoid with centre $-\frac{P}{2} \bfI$, hence strictly convex with $C^\infty$ boundary, and the outer unit normal $\partial \K(P) \rightarrow S^2$ is bijective.
    \end{enumerate}
\end{lemma}

\begin{assumption}\label{ass}
    \begin{enumerate}
        \item[(A1)] Let $\Omega \subset \mathbb R^2$ be a bounded Lipschitz domain and $\partial \Omega = \bar \Gamma_D \cup \bar \Gamma_N$ with relatively open and disjoint $\Gamma_D, \Gamma_N$ with $\mathcal{H}^1(\partial \Omega \setminus (\Gamma_D \cup \Gamma_N))=0$. The case $\Gamma_D = \emptyset$ is allowed.
        \item[(A2)] Let $m \in L^\infty(\Omega)$, $m\geq 0$, $c_o>0$ constant and $f_c \in L^\infty(\Omega)$.
        \item[(A3)] Let $P \in C(\bar \Omega)$ and $P\geq0$, in particular $P=0$
        on subsets (ice boundary and open water) is allowed.
        \item[(A4)] Let $\sqrt{m} \bfu^{n-1} \in L^2(\Omega)^2$, $\bfu_o \in L^3(\Omega)^2$ and $\bff \in L^{3/2}(\Omega)^2$.
    \end{enumerate}
\end{assumption}

\begin{lemma}\label{lem:prop_D}
    For $P\geq 0$ and $\bfxi \in \mathbb S$ we have
    \begin{align}
        D(P,\bfxi) := \sup_{\bftau\in\K(P)} \bftau : \bfxi = \frac{P}{2}(\Delta(\bfxi)-\xi_I) \ .
        \label{eq:def_D}
    \end{align}
    In particular $D(P,\cdot)$ is convex, homogeneous of degree 1, non-negative,
    Lipschitz continuous with constant $c_\K P$, and we have
    \begin{align}
        D(P,\bfxi) = 0 \Leftrightarrow \bfxi = \alpha \bfI \ \text{for some}\ \alpha \geq 0 \label{eq:dissipationfree}
    \end{align}
    for all $P>0$. Thus isotropic expansion is dissipation-free. Positive
divergence alone does not imply zero dissipation. This property does not
exclude positive principal stresses elsewhere on the yield ellipse.
Additionally the pointwise estimate
    \begin{align}
        D(P, \bfxi) \geq P(\xi_I)^- \label{eq:negativ_part_D}
    \end{align}
    holds, where $(\cdot)^-$ is the negative part.
\end{lemma}

    \begin{align}
        \langle \bftau\nu, \bfw \rangle := \int_\Omega \bftau : \varepsilon(\bfw) \dd x + \int_\Omega \bfw \cdot \div \bftau \dd x, \quad\text{for all }\bfw\in W^{1,3}(\Omega)^2\ . \label{eq:normaltrace}
    \end{align}

\begin{lemma}\label{lem:a2}
    Let $\bftau \in \Sigma$ and $\bfw \in W^{1,3}(\Omega)^2$ with $\bfw =0$ a.e.\ in $U\cap \Omega$ for an open set $U\subseteq \mathbb R^2$
    with $\bar\Gamma_D \subset U$. Then $\langle \bftau\nu, \bfw \rangle =0$.
\end{lemma}

\begin{lemma}\label{lem:a3}
    Under Assumption \ref{ass}, for every $\bftau\in\Sigma$ and $\bfv \in W^{1,3}(\Omega)^2$
    \begin{align*}
        - \langle\bftau\nu,\bfv\rangle \leq \int_{\Gamma_D} D(P, -(\bfv\odot \nu)) \dd \mathcal{H}^1
    \end{align*}
    where $\bfv$ on $\Gamma_D$ is evaluated in the sense of traces.
\end{lemma}

\begin{proof}
    Step 1:
    Let $\eta>0$ and $\zeta_\eta(x) := \psi (\dist(x,\bar\Gamma_D)/\eta)$ with $\psi \in \mathrm{Lip}(\mathbb R)$
    and $\psi = 1$ in $(-\infty,1/2]$, $\psi=0$ in $[1,\infty)$ and $0\leq \psi \leq 1$.
    Then $(1-\zeta_\eta)\bfv =0$ in $U_{\eta/2}(\bar\Gamma_D)\cap \Omega$. By Lemma \ref{lem:a2} we have
    $\langle \bftau \nu, \bfv\rangle = \langle \bftau\nu,\zeta_\eta \bfv\rangle$. We will prove that
    \begin{align}
        -\langle\bftau\nu,\bfw\rangle \leq \int_{\partial\Omega} D(P,-(\bfw\odot \nu)) \dd \mathcal{H}^1 \label{eq:step1}
    \end{align}
    for $\bfw:=\zeta_\eta\bfv \in W^{1,3}(\Omega)^2$, and we take the limit as $\eta\rightarrow0$
    in Step 6.

    Step 2:
    Assumption \ref{ass} (A1) implies that there exist finitely many open cylinders $Z_1,\dots,Z_N \subset \mathbb R^2$
    which cover $\partial\Omega$, such that (after a rigid motion)
    \begin{align*}
        \Omega\cap Z_j = \{y=(y_1,y_2)\in Z_j : y_2>\gamma_j(y_1)\}, \quad \partial\Omega\cap Z_j
        =\{y:y_2=\gamma_j(y_1)\}\cap Z_j
    \end{align*}
    with $L$-Lipschitz continuous functions $\gamma_j$. The outer normal and the surface element
    in these coordinates are, for $\mathcal{L}^1$ a.e.\ $y_1$,
    \begin{align*}
        \nu(y_1) = \frac{(\gamma_j'(y_1),-1)}{\sqrt{1+\gamma'_j(y_1)^2}}, \quad \dd \mathcal{H}^1 = \sqrt{1+\gamma'_j(y_1)^2} \dd y_1 \ .
    \end{align*}
    Both are invariant under vertical translations. Choose $Z_0\Subset \Omega$ with
    $\bar\Omega \subset \cup_{j=0}^N Z_j$ and a subordinate smooth partition of unity $(\theta_j)_{j=0}^N$,
    $\theta_j\in C^\infty_c(Z_j)$ and $\sum_{j=0}^{N}\theta_j =1$ on $\bar\Omega$.

    Since \eqref{eq:normaltrace} is linear in $\bfw$ and
    $\theta_j\bfw \in W^{1,3}(\Omega)$, we notice that \eqref{eq:step1} can be shown term by term.
    The term for $j=0$ vanishes, since $\theta_0\bfw\in W^{1,3}$ has compact support
    in $\Omega$. As in the proof of Lemma \ref{lem:a2}, the mollified functions have zero trace on $\partial\Omega$ and the limit leads to
    $\langle \bftau\nu, \theta_0\bfw\rangle =0$.

    From now on let $j\geq 1$ be fixed, $\gamma:= \gamma_j$, and redefine $\bfw:= \theta_j\zeta_\eta\bfv$. Choose an open interval $B' \Subset \{y_1:(y_1,y_2)\in Z_j\ \text{for a}\ y_2\}$, a height $M$ and $t_0>0$ with
    \begin{align*}
        \supp \bfw \cap \bar\Omega \subset \{y : y_1\in B', \gamma(y_1)\leq y_2 < M-t_0\} \Subset Z_j \ .
    \end{align*}
    Step 3: For $t \in (0,t_0)$ we define
    \begin{align*}
        D_t := \{y : y_1\in B', \gamma(y_1)+t <y_2<M\}, \quad S_t =\{y : y_1\in B', y_2 = \gamma(y_1)+t\} \ .
    \end{align*}
    $D_t$ is a Lipschitz domain and $S_t$ its lower boundary. We notice that $\bfw$ vanishes for $y_1\in\partial B'$ and $y_2=M$. For $y\in D_t$ a simple calculation shows that $\dist(y,\partial\Omega)\geq c_L t$ for $c_L := (1+L^2)^{-1/2}$ and
    $L$ the Lipschitz constant of $\gamma$.

    Let $\bftau_\rho := \bftau * \varphi_\rho$ with $\varphi_\rho$ the standard mollifier and $\rho < c_L t/2$. On a neighbourhood of $\bar D_t$, $\bftau_\rho$ is
    smooth and we have
     \begin{align*}
        \div \bftau_\rho = (\div\bftau) * \varphi_\rho \ ,
     \end{align*}
     and with the outer normal $\nu(y_1)$ of $D_t$ on $S_t$ (identical to the one of $\Omega$, by
     translation invariance)
     \begin{align}
        \int_{D_t} \bftau_\rho : \varepsilon(\bfw) + \bfw \cdot \div\bftau_\rho \dd y =
        \int_{S_t} \bfw \cdot(\bftau_\rho \nu) \dd \mathcal{H}^1 \ . \label{eq:trans}
     \end{align}
     Step 4:
     The function $T(y_1, s):= (y_1, \gamma(y_1)+s)$ is bi-Lipschitz from $B'\times (0,t_0)$ to
     \begin{align*}
        C := T(B'\times (0,t_0)) = \{y : y_1\in B', 0<y_2-\gamma(y_1)<t_0\} = \cup_{0<t<t_0} S_t
     \end{align*}
     and measure preserving ($\det \mathrm{D}T=1$). In the following we will apply
     a Fubini argument to select admissible slices.
     Let $N_1 \subset \Omega$ be the $\mathcal{L}^2$ null set of non-Lebesgue points of $\bftau$,
     $N_2$ the null set where $\bftau(x)\in \K(P(x))$ fails and $N_3$ the null set where $\bfw \circ T$ is not absolutely continuous on lines (ACL) in $s$-direction. With Fubini we have for $\mathcal{L}^1$ a.e.\ $t\in(0,t_0)$ that $\mathcal{H}^1$ a.e.\ points of $S_t$ are Lebesgue points of $\bftau$ and admissible.

    For such a $t$ we take the limit as $\rho \rightarrow 0$ in \eqref{eq:trans}.
    For the right-hand side, $\bftau_\rho \rightarrow \bftau$ pointwise $\mathcal{H}^1$ a.e.\ on $S_t$ with majorant $\Vert \bftau\Vert_{L^\infty}|\bfw_{|S_t}| \in L^1(S_t)$.
    With the pointwise bound of the support function (Lemma \ref{lem:prop_D}) we have
    \begin{align*}
        \bfw \cdot (\bftau\nu) = \bftau : (\bfw\odot \nu) \geq - D(P, -(\bfw\odot \nu))
    \end{align*}
    $\mathcal{H}^1$ a.e.\ on $S_t$. For a.e.\ $t\in(0,t_0)$ we therefore have
    \begin{align}
        -\int_{D_t} \bftau : \varepsilon(\bfw) + \bfw\cdot\div\bftau \dd y
        \leq \int_{S_t} D(P, -(\bfw\odot \nu)) \dd \mathcal{H}^1 \ . \label{eq:slice}
    \end{align}
    Step 5:
    For the left-hand side, the monotone exhaustion $D_t \nearrow D_0$ with
    $L^1$ integrand leads to convergence to $-\langle \bftau\nu, \bfw\rangle$, since
    the support of $\bfw$ makes the $D_0$ integral equal to the $\Omega$ integral.
    For the right-hand side we use the translation invariance of $\nu$ and $\dd \mathcal{H}^1$:
    \begin{align*}
        &\int_{S_t} D(P, -(\bfw\odot \nu)) \dd \mathcal{H}^1\\
        &\quad=\int_{B'} D(P(T(y_1,t)), -(\bfw (T(y_1,t))\odot \nu(y_1))) \sqrt{1+\gamma'(y_1)^2} \dd y_1 \ .
    \end{align*}
    By ACL we have for a.e.\ $y_1$
    \begin{align*}
        |\bfw (T(y_1,t)) - \bfw (T(y_1,0^+))| \leq \int_{0}^{t} |\partial_s(\bfw\circ T)(y_1,s)| \dd s,
    \end{align*}
    and integration over $y_1$ and H\"older lead to
    \begin{align*}
        \Vert \bfw \circ T(\cdot, t) - \bfw \circ T(\cdot, 0^+)\Vert_{L^1(B')}
        \leq  C t^{2/3} \Vert \nabla(\bfw \circ T)\Vert_{L^3} \rightarrow 0 \ .
    \end{align*}
    The ACL boundary values $\bfw \circ T(\cdot, 0^{+})$ coincide $\mathcal{L}^{1}$ a.e.\ with the trace $\bfw_{|\partial\Omega}$.
    Also $|P(T(y_1,t))-P(T(y_1,0))| \leq \omega_P(t)$, where $\omega_P$ is the modulus of continuity of $P$.
    The Lipschitz continuity from Lemma \ref{lema:prop_K} and Lemma \ref{lem:prop_D} lead to
    \begin{align*}
        |D(P_1,\bfxi_1) - D(P_2,\bfxi_2)| \leq c_\K |P_1-P_2| |\bfxi_1|_F + c_\K \Vert P\Vert_{L^\infty}
        |\bfxi_1-\bfxi_2|_F \ .
    \end{align*}
    This and $|a\odot b|_F \leq |a||b|$ lead to the convergence of the right-hand side of \eqref{eq:slice}
    to $\int_{\partial\Omega} D(P,-(\bfw\odot \nu)) \dd \mathcal{H}^1$ along a sequence of admissible $t$. This proves \eqref{eq:step1}, after summation over $j$.

    Step 6:
    Returning to Step 1 and \eqref{eq:step1}, we have with $\bfw = \zeta_\eta\bfv$ and the positive homogeneity
    of degree 1 of $D(P,\cdot)$ (Lemma \ref{lem:prop_D}) and $0\leq\zeta_\eta\in\mathrm{Lip}(\bar\Omega)$
    \begin{align*}
        -\langle \bftau\nu, \bfv\rangle = - \langle \bftau\nu, \zeta_\eta \bfv\rangle
        \leq \int_{\partial\Omega} \zeta_\eta D(P, -(\bfv\odot \nu)) \dd \mathcal{H}^1 \ .
    \end{align*}
    As $\eta\rightarrow 0$, $\zeta_\eta \rightarrow \mathbf 1_{\bar\Gamma_D}$ pointwise on $\partial\Omega$.
    The integrands are bounded by the $\mathcal{H}^{1}$-integrable function
    $c_\K \Vert P\Vert_{L^{\infty}}|\bfv_{|\partial\Omega}|$.
    Dominated convergence and $\mathcal{H}^{1}(\bar\Gamma_D\setminus\Gamma_D) =0$ conclude the proof.
\end{proof}

\end{document}
